\begin{textsample}{POS Dim 1 – vem\_ed\_23 – Score 20.00 – vem\_ed\_23/t123.txt}  \label{ex:f1_pos_139}
Quem \textbf{é} Quem Cathy Lee T.
Arcuino \textbf{é} diretora executiva de Global Engagement (em tradução livre, engajamento para a cidadania global) na Wilkes University (EUA).
Logo após sua graduação, em 1999, atuou como voluntária do Peace Corps em uma \textbf{pequena} vila no sul do Cazaquistão, próximo à fronteira uzbeque.
Ensinou inglês no Japão, na Tailândia e na Polônia, trabalhou numa ONG no Quirguistão e participou do \textbf{programa} Fulbright International Educator Administrator no Japão.
Mestre e doutora em Educação (respectivamente, pela Framingham State College e pela Colorado State University), \textbf{é} formada pela NAFSA Academy (Associação de Educadores Internacionais).
Conte como suas origens influenciaram sua visão de cidadania global.
Meus pais \textbf{são} filipinos \textbf{que} emigraram para os Estados Unidos.
Nasci em Nova York e quando tinha 4 anos nos mudamos para a Califórnia.
Meus pais me ensinaram \textbf{que} eu teria \textbf{que} trabalhar o dobro, não só porque \textbf{sou} uma mulher nos Estados Unidos, mas por \textbf{ser} asiática.
Tenho doutorado em Educação e meus pais ficaram felizes porque queriam \textbf{que} eu \textbf{fosse} “doutora” médica, mas descobriram \textbf{que} me \textbf{tornara} Dra.
Arcuino.
Na juventude, visitava minha família nas Filipinas e \textbf{observava} os recursos educacionais \textbf{que} tive e meus primos não tiveram lá.
E me perguntava: e se meus pais não tivessem emigrado?
Como percebo \textbf{que} esse \textbf{pequeno} movimento \textbf{que} meus pais fizeram mudou completamente minha vida, quero transformar outras vidas por meio da internacionalização.
Fale sobre sua experiência no Peace Corps.
Quando ingressei como voluntária no Peace Corps, minha \textbf{expectativa} \textbf{era} ir para as Filipinas.
Só \textbf{que} me disseram \textbf{que} teria de aguardar um ano e não queria esperar tanto, então \textbf{fui} parar no Cazaquistão.
Para \textbf{ser} honesta, tinha 21 anos e não fazia ideia de onde ficava o país. \textbf{Foi} uma das melhores decisões \textbf{que} fiz porque conheci um país e uma cultura lindas.
Também \textbf{aprendi} muito sobre mim mesma.
Não falava nada de russo, só \textbf{sabia} \textbf{que} “niet” significa “não”.
Naquela época o russo \textbf{era} a língua principal, devido à influência soviética, especialmente no sul do país, onde vivi, próximo à fronteira uzbeque.
Como \textbf{foi} sua experiência como professora de inglês no Japão? \textbf{Foi} um grande um choque cultural.
Logo depois de viver numa \textbf{pequena} vila cazaque, \textbf{fui} para Osaka, Japão.
Também \textbf{foi} muito diferente porque naquela época [ 2001 ] o Japão estava muito avançado no ensino-aprendizagem online.
Eu vivia em Osaka, mas meus alunos \textbf{eram} de \textbf{todo} o país. \textbf{Era} exatamente como estamos nos encontrando agora.
Não tinha mais do \textbf{que} três estudantes online e eu ensinava inglês.
Nunca poderia imaginar \textbf{que} vinte e poucos anos depois isso \textbf{seria} \textbf{parte} do mundo, porque naquela época isso já estava acontecendo no Japão.
No Japão se trabalha muitas horas por dia, então para fazer cursos extracurriculares ou para aprimorar \textbf{habilidades} \textbf{linguísticas}, \textbf{é} necessário fazer isso após o trabalho e o transporte de volta para \textbf{casa}.
Então meu turno \textbf{começava} por volta das 22h30 e terminava 6h30, 7h da manhã.
Muitos estudavam inglês na madrugada.
Esse desejo de \textbf{aprender} \textbf{era} fascinante.
Tanto a experiência no Cazaquistão \textbf{quanto} em Osaka me fizeram ver \textbf{que} estou no campo certo, porque sinto \textbf{que} realmente estou contribuindo com os objetivos educacionais das \textbf{pessoas}.
Depois disso, vivi dois anos na Tailândia ensinando inglês, e outros dois anos na Polônia.
Enfim, senti saudades da Ásia Central e me mudei para o Quirguistão, onde ajudei estudantes quirguizes a estudar nos Estados Unidos.
Finalmente voltei para os Estados Unidos.
Visitando universidades ao redor do mundo, o \textbf{que} poderia dizer sobre Internacionalização e Mobilidade?
Diria \textbf{que} \textbf{é} fascinante ver o crescimento e o interesse [ sobre esses temas ].
Primeiro, devido ao \textbf{desenvolvimento} da tecnologia, crescendo globalmente na ponta dos nossos dedos.
Isso leva a geração mais nova a se interessar e \textbf{aprender} rapidamente, instantaneamente, sobre o \textbf{que} acontece no mundo.
Sinto \textbf{que} a mobilidade \textbf{é} mais rápida, além de mais acessível, porque mesmo \textbf{que} você não consiga viajar fisicamente, você pode \textbf{aprender} sobre outros países assistindo \textbf{conteúdos} online ou fazendo \textbf{conexões} e mantendo melhor os contatos.
Veja como estamos falando agora, eu na Pensilvânia e vocês aí em São Paulo.
Outra \textbf{questão} \textbf{é} \textbf{que} existe mais \textbf{conhecimento} \textbf{sendo} \textbf{compartilhado}.
No tempo em \textbf{que} vivi no Cazaquistão, a não \textbf{ser} \textbf{que} você conhecesse alguém de lá, você não \textbf{saberia} sobre o país a não \textbf{ser} a \textbf{partir} das notícias.
Agora, há outros caminhos.
Por exemplo, a música, \textbf{que} se \textbf{tornou} mais global.
Você vê de fato o \textbf{impacto} da mobilidade global, antes \textbf{era} um presente \textbf{que} apenas poucas \textbf{pessoas} conseguiam.
Agora \textbf{é} mais acessível a \textbf{todos}.
O \textbf{que} considero mais importante \textbf{é} o interesse de \textbf{pessoas} de mente aberta \textbf{que} desejam estar juntas.
Ver paisagens online desperta mais vontade de viajar.
Mas, ainda assim, existe a limitação financeira para viajar fisicamente.
Existe muito medo de viajar a \textbf{partir} de uma \textbf{pequena} cidade para outro lugar do mundo.
Qual sua experiência \textbf{profissional} mais desafiadora e gratificante?
Culturalmente, em \textbf{todos} os lugares em \textbf{que} estive houve uma revelação interessante, mas também um desafio.
Trabalhando com tantas culturas diferentes, há complexidades tão diversas.
Estou constantemente trocando de código.
O \textbf{que} funciona com os brasileiros pode não funcionar com nigerianos ou espanhóis.
Diria \textbf{que} a maioria dos meus desafios \textbf{profissionais} se relaciona com me assegurar de \textbf{que} não esteja \textbf{sendo} ofensiva e \textbf{reconhecer} \textbf{que} a \textbf{forma} como vou me comunicar com um país pode não \textbf{ser} o mesmo método \textbf{que} preciso \textbf{utilizar} com outro.
Gostaria \textbf{que} você comentasse a noção de “global engagement”, \textbf{que} não \textbf{é} muito conhecida aqui no Brasil.
Do meu \textbf{ponto} de vista às vezes a \textbf{questão} \textbf{é} interpretada assim: “\textbf{quantos} estrangeiros entram na universidade e \textbf{quantos} estudantes da instituição saem”.
Mas quando se está trabalhando pelo engajamento global, quais \textbf{são} os outros caminhos?
Adoro as parcerias COIL [ Collaborative Online International Learning ], porque \textbf{reconhecemos} \textbf{que} nem \textbf{todos} terão oportunidades de ir para outro país.
Mas poderão \textbf{aprender} muito virtualmente, por exemplo, sobre as comidas do outro país.
Outra maneira de engajar globalmente \textbf{é} apoiar professores, por exemplo, pesquisando em dois países e áreas diferentes.
Uma outra \textbf{forma} \textbf{é} engajar a comunidade \textbf{local} e despertar a curiosidade sobre compreender outras culturas.
Esse \textbf{é} um bom começo.
Por outro lado, existe muito medo do desconhecido.
Respondendo à sua \textbf{questão}, o conceito de \textbf{ser} engajado globalmente se relaciona com algo feito fisicamente através da sua instituição, mas também por meio de sua curiosidade, de maneira informal. \textbf{Que} \textbf{futuro} você vislumbra para os projetos COIL na sua instituição e em \textbf{geral} para os próximos cinco anos?
Uma das \textbf{coisas} \textbf{que} vejo \textbf{é} o aumento nas colaborações COIL.
Meu objetivo \textbf{é} \textbf{criar} muitos \textbf{pontos} de contatos entre estudantes e professores por meio de iniciativas COIL.
Também tenho a visão de triangular as colaborações em projetos com mais de dois países.

% matched lemmas: aprender, casa, coisa, começar, compartilhar, conexão, conhecimento, conteúdo, criar, desenvolvimento, expectativa, forma, futuro, geral, habilidade, impacto, linguístico, local, observar, parte, partir, pequeno, pessoa, ponto, profissional, programa, quanto, que, questão, reconhecer, saber, ser, todo, tornar, utilizar
\end{textsample}
